\Chapter{93}

\Verse{1}
{
    \zR{却说冯紫英去后，贾政叫门上的人来吩咐道：“今儿临安伯那里来请吃酒，知 道是什么事？”}
}
{
    \eR{[English source not present in the checked-in Bencraft/Joly files.]}
}
{
}

\Verse{2}
{
    \zR{门上的人道：“奴才曾问过，并没有什么喜庆事，不过南安王府里 到了一班小戏子，都说是个名班，伯爷高兴，唱两天戏请相好的老爷们瞧瞧，热闹 热闹。}
}
{
    \eR{[English source not present in the checked-in Bencraft/Joly files.]}
}
{
}

\Verse{3}
{
    \zR{大约不用送礼的。”}
}
{
    \eR{[English source not present in the checked-in Bencraft/Joly files.]}
}
{
}

\Verse{4}
{
    \zR{说着，贾赦过来问道：“明儿二老爷去不去？”}
}
{
    \eR{[English source not present in the checked-in Bencraft/Joly files.]}
}
{
}

\Verse{5}
{
    \zR{贾政道： “承他亲热，怎么好不去的。”}
}
{
    \eR{[English source not present in the checked-in Bencraft/Joly files.]}
}
{
}

\Verse{6}
{
    \zR{说着，门上进来回道：“衙门里书办来请老爷明日 上衙门。}
}
{
    \eR{[English source not present in the checked-in Bencraft/Joly files.]}
}
{
}

\Verse{7}
{
    \zR{有堂派的事，必得早些去。”}
}
{
    \eR{[English source not present in the checked-in Bencraft/Joly files.]}
}
{
}

\Verse{8}
{
    \zR{贾政道：“知道了。”}
}
{
    \eR{[English source not present in the checked-in Bencraft/Joly files.]}
}
{
}

\Verse{9}
{
    \zR{说着，只见两个管屯 里地租子的家人走来，请了安磕了头旁边站着。}
}
{
    \eR{[English source not present in the checked-in Bencraft/Joly files.]}
}
{
}

\Verse{10}
{
    \zR{贾政道：“你们是郝家庄的？”}
}
{
    \eR{[English source not present in the checked-in Bencraft/Joly files.]}
}
{
}

\Verse{11}
{
    \zR{两 个答应了一声。}
}
{
    \eR{[English source not present in the checked-in Bencraft/Joly files.]}
}
{
}

\Verse{12}
{
    \zR{贾政也不往下问，竟与贾赦各自说了一回话儿散了。}
}
{
    \eR{[English source not present in the checked-in Bencraft/Joly files.]}
}
{
}

\Verse{13}
{
    \zR{家人等秉着手灯送过贾赦去，这里贾琏便叫那管租的人道：“说你的。”}
}
{
    \eR{[English source not present in the checked-in Bencraft/Joly files.]}
}
{
}

\Verse{14}
{
    \zR{那人 说道：“十月里的租子，奴才已经赶上来了。}
}
{
    \eR{[English source not present in the checked-in Bencraft/Joly files.]}
}
{
}

\Verse{15}
{
    \zR{原是明儿可到，谁知京外拿车，把车 上的东西不由分说都掀在地下。}
}
{
    \eR{[English source not present in the checked-in Bencraft/Joly files.]}
}
{
}

\Verse{16}
{
    \zR{奴才告诉他，说是府里收租子的车，不是买卖车， 他更不管这些。}
}
{
    \eR{[English source not present in the checked-in Bencraft/Joly files.]}
}
{
}

\Verse{17}
{
    \zR{奴才叫车夫只管拉着走，几个衙役就把车夫混打了一顿，硬扯了两 辆车去了。}
}
{
    \eR{[English source not present in the checked-in Bencraft/Joly files.]}
}
{
}

\Verse{18}
{
    \zR{奴才所以先来回报。}
}
{
    \eR{[English source not present in the checked-in Bencraft/Joly files.]}
}
{
}

\Verse{19}
{
    \zR{求爷打发个人到衙门里去要了来才好。}
}
{
    \eR{[English source not present in the checked-in Bencraft/Joly files.]}
}
{
}

\Verse{20}
{
    \zR{再者，也整 治整治这些无法无天的差役才好。}
}
{
    \eR{[English source not present in the checked-in Bencraft/Joly files.]}
}
{
}

\Verse{21}
{
    \zR{爷还不知道呢：更可怜的是那买卖车，客商的东 西全不顾，掀下来赶着就走。}
}
{
    \eR{[English source not present in the checked-in Bencraft/Joly files.]}
}
{
}

\Verse{22}
{
    \zR{那些赶车的但说句话，打的头破血出的。”}
}
{
    \eR{[English source not present in the checked-in Bencraft/Joly files.]}
}
{
}

\Verse{23}
{
    \zR{贾琏听了， 骂道：“这个还了得！”}
}
{
    \eR{[English source not present in the checked-in Bencraft/Joly files.]}
}
{
}

\Verse{24}
{
    \zR{立刻写了一个帖儿，叫家人：“拿去向拿车的衙门里要车 去，并车上东西，若少了一件是不依的。}
}
{
    \eR{[English source not present in the checked-in Bencraft/Joly files.]}
}
{
}

\Verse{25}
{
    \zR{快叫周瑞。”}
}
{
    \eR{[English source not present in the checked-in Bencraft/Joly files.]}
}
{
}

\Verse{26}
{
    \zR{周瑞不在家，又叫旺儿。}
}
{
    \eR{[English source not present in the checked-in Bencraft/Joly files.]}
}
{
}

\Verse{27}
{
    \zR{旺 儿晌午出去了，还没有回来。}
}
{
    \eR{[English source not present in the checked-in Bencraft/Joly files.]}
}
{
}

\Verse{28}
{
    \zR{贾琏道：“这些忘八日的，一个都不在家!他们成年 家吃粮不管事！”}
}
{
    \eR{[English source not present in the checked-in Bencraft/Joly files.]}
}
{
}

\Verse{29}
{
    \zR{因吩咐小厮们：“快给我找去！”}
}
{
    \eR{[English source not present in the checked-in Bencraft/Joly files.]}
}
{
}

\Verse{30}
{
    \zR{说着，也回到自己屋里睡下， 不提。}
}
{
    \eR{[English source not present in the checked-in Bencraft/Joly files.]}
}
{
}

\Verse{31}
{
    \zR{且说临安伯第二天又打发人来请。}
}
{
    \eR{[English source not present in the checked-in Bencraft/Joly files.]}
}
{
}

\Verse{32}
{
    \zR{贾政告诉贾赦道：“我是衙门里有事。}
}
{
    \eR{[English source not present in the checked-in Bencraft/Joly files.]}
}
{
}

\Verse{33}
{
    \zR{琏儿 要在家等候拿车的事情，也不能去。}
}
{
    \eR{[English source not present in the checked-in Bencraft/Joly files.]}
}
{
}

\Verse{34}
{
    \zR{倒是大老爷带着宝玉应酬一天也罢了。”}
}
{
    \eR{[English source not present in the checked-in Bencraft/Joly files.]}
}
{
}

\Verse{35}
{
    \zR{贾赦 点头道：“也使得。”}
}
{
    \eR{[English source not present in the checked-in Bencraft/Joly files.]}
}
{
}

\Verse{36}
{
    \zR{贾政遣人去叫宝玉，说：“今儿跟大爷到临安伯那里听戏去。”}
}
{
    \eR{[English source not present in the checked-in Bencraft/Joly files.]}
}
{
}

\Verse{37}
{
    \zR{宝玉喜欢的了不得，便换上衣服，带了焙茗、扫红、锄药三个小子，出来见了贾赦， 请了安，上了车，来到临安伯府里。}
}
{
    \eR{[English source not present in the checked-in Bencraft/Joly files.]}
}
{
}

\Verse{38}
{
    \zR{门上人回进去，一会子出来说：“老爷请。”}
}
{
    \eR{[English source not present in the checked-in Bencraft/Joly files.]}
}
{
}

\Verse{39}
{
    \zR{于是贾赦带着宝玉走入院内，只见宾客喧阗。}
}
{
    \eR{[English source not present in the checked-in Bencraft/Joly files.]}
}
{
}

\Verse{40}
{
    \zR{贾赦宝玉见了临安伯，又与众宾客都 见过了礼，大家坐着，说笑了一回。}
}
{
    \eR{[English source not present in the checked-in Bencraft/Joly files.]}
}
{
}

\Verse{41}
{
    \zR{只见一个掌班拿着一本戏单，一个牙笏，向上 打了一个千儿，说道：“求各位老爷赏戏。”}
}
{
    \eR{[English source not present in the checked-in Bencraft/Joly files.]}
}
{
}

\Verse{42}
{
    \zR{先从尊位点起，挨至贾赦，也点了一 出。}
}
{
    \eR{[English source not present in the checked-in Bencraft/Joly files.]}
}
{
}

\Verse{43}
{
    \zR{那人回头见了宝玉，便不向别处去，竟抢步上来，打个千儿道：“求二爷赏两 出。”}
}
{
    \eR{[English source not present in the checked-in Bencraft/Joly files.]}
}
{
}

\Verse{44}
{
    \zR{宝玉一见那人，面如傅粉，唇若涂朱，鲜润如出水芙渠，飘扬似临风玉树： 原来不是别人，就是蒋玉函。}
}
{
    \eR{[English source not present in the checked-in Bencraft/Joly files.]}
}
{
}

\Verse{45}
{
    \zR{前日听得他带了小戏儿进京，也没有到自己那里；此 时见了，又不好站起来，只得笑道：“你多早晚来的？”}
}
{
    \eR{[English source not present in the checked-in Bencraft/Joly files.]}
}
{
}

\Verse{46}
{
    \zR{蒋玉函把眼往左右一溜， 悄悄的笑道：“怎么二爷不知道么？”}
}
{
    \eR{[English source not present in the checked-in Bencraft/Joly files.]}
}
{
}

\Verse{47}
{
    \zR{宝玉因众人在坐，也难说话，只得乱点了一 出。}
}
{
    \eR{[English source not present in the checked-in Bencraft/Joly files.]}
}
{
}

\Verse{48}
{
    \zR{蒋玉函去了，便有几个议论道：“此人是谁？”}
}
{
    \eR{[English source not present in the checked-in Bencraft/Joly files.]}
}
{
}

\Verse{49}
{
    \zR{有的说：“他向来是唱小旦的， 如今不肯唱小旦，年纪也大了，就在府里掌班。}
}
{
    \eR{[English source not present in the checked-in Bencraft/Joly files.]}
}
{
}

\Verse{50}
{
    \zR{头里也改过小生。}
}
{
    \eR{[English source not present in the checked-in Bencraft/Joly files.]}
}
{
}

\Verse{51}
{
    \zR{他也攒了好几个 钱，家里已经有两三个铺子，只是不肯放下本业，原旧领班。”}
}
{
    \eR{[English source not present in the checked-in Bencraft/Joly files.]}
}
{
}

\Verse{52}
{
    \zR{有的说：“想必成 了家了。”}
}
{
    \eR{[English source not present in the checked-in Bencraft/Joly files.]}
}
{
}

\Verse{53}
{
    \zR{有的说：“亲还没有定。}
}
{
    \eR{[English source not present in the checked-in Bencraft/Joly files.]}
}
{
}

\Verse{54}
{
    \zR{他倒拿定一个主意，说是人生婚配关系一生一 世的事，不是混闹得的，不论尊卑贵贱，总要配的上他的才能。}
}
{
    \eR{[English source not present in the checked-in Bencraft/Joly files.]}
}
{
}

\Verse{55}
{
    \zR{所以到如今还并没 娶亲。”}
}
{
    \eR{[English source not present in the checked-in Bencraft/Joly files.]}
}
{
}

\Verse{56}
{
    \zR{宝玉暗忖度道：“不知日后谁家的女孩儿嫁他？}
}
{
    \eR{[English source not present in the checked-in Bencraft/Joly files.]}
}
{
}

\Verse{57}
{
    \zR{要嫁着这么样的人才儿， 也算是不辜负了。”}
}
{
    \eR{[English source not present in the checked-in Bencraft/Joly files.]}
}
{
}

\Verse{58}
{
    \zR{那时开了戏，也有昆腔，也有高腔，也有弋腔、平腔，热闹非常。}
}
{
    \eR{[English source not present in the checked-in Bencraft/Joly files.]}
}
{
}

\Verse{59}
{
    \zR{到了晌午， 便摆开桌子吃酒。}
}
{
    \eR{[English source not present in the checked-in Bencraft/Joly files.]}
}
{
}

\Verse{60}
{
    \zR{又看了一回，贾赦便欲起身。}
}
{
    \eR{[English source not present in the checked-in Bencraft/Joly files.]}
}
{
}

\Verse{61}
{
    \zR{临安伯过来留道：“天色尚早。}
}
{
    \eR{[English source not present in the checked-in Bencraft/Joly files.]}
}
{
}

\Verse{62}
{
    \zR{听 见说琪官儿还有一出《占花魁》，他们顶好的首戏。”}
}
{
    \eR{[English source not present in the checked-in Bencraft/Joly files.]}
}
{
}

\Verse{63}
{
    \zR{宝玉听了，巴不得贾赦不走。}
}
{
    \eR{[English source not present in the checked-in Bencraft/Joly files.]}
}
{
}

\Verse{64}
{
    \zR{于是贾赦又坐了一会。}
}
{
    \eR{[English source not present in the checked-in Bencraft/Joly files.]}
}
{
}

\Verse{65}
{
    \zR{果然蒋玉函扮了秦小官，伏侍花魁醉后神情，把那一种怜香 惜玉的意思，做得极情尽致。}
}
{
    \eR{[English source not present in the checked-in Bencraft/Joly files.]}
}
{
}

\Verse{66}
{
    \zR{以后对饮对唱，缠绵缱绻。}
}
{
    \eR{[English source not present in the checked-in Bencraft/Joly files.]}
}
{
}

\Verse{67}
{
    \zR{宝玉这时不看花魁，只把 两支眼睛独射在秦小官身上。}
}
{
    \eR{[English source not present in the checked-in Bencraft/Joly files.]}
}
{
}

\Verse{68}
{
    \zR{更加蒋玉函声音响亮，口齿清楚，按腔落板，宝玉的 神魂都唱的飘荡了。}
}
{
    \eR{[English source not present in the checked-in Bencraft/Joly files.]}
}
{
}

\Verse{69}
{
    \zR{直等这出戏煞场后，更知蒋玉函极是情种，非寻常脚色可比。}
}
{
    \eR{[English source not present in the checked-in Bencraft/Joly files.]}
}
{
}

\Verse{70}
{
    \zR{因想着：“《乐记》上说的是：‘情动于中，故形于声；声成文，谓之音。’}
}
{
    \eR{[English source not present in the checked-in Bencraft/Joly files.]}
}
{
}

\Verse{71}
{
    \zR{所以 知声，知音，知乐，有许多讲究。}
}
{
    \eR{[English source not present in the checked-in Bencraft/Joly files.]}
}
{
}

\Verse{72}
{
    \zR{声音之原，不可不察。}
}
{
    \eR{[English source not present in the checked-in Bencraft/Joly files.]}
}
{
}

\Verse{73}
{
    \zR{诗词一道，但能传情，不 能入骨，自后想要讲究讲究音律。”}
}
{
    \eR{[English source not present in the checked-in Bencraft/Joly files.]}
}
{
}

\Verse{74}
{
    \zR{宝玉想出了神，忽见贾赦起身，主人不及相留。}
}
{
    \eR{[English source not present in the checked-in Bencraft/Joly files.]}
}
{
}

\Verse{75}
{
    \zR{宝玉没法，只得跟了回来。}
}
{
    \eR{[English source not present in the checked-in Bencraft/Joly files.]}
}
{
}

\Verse{76}
{
    \zR{到了家中，贾赦自回那边去了。}
}
{
    \eR{[English source not present in the checked-in Bencraft/Joly files.]}
}
{
}

\Verse{77}
{
    \zR{宝玉来见贾政。}
}
{
    \eR{[English source not present in the checked-in Bencraft/Joly files.]}
}
{
}

\Verse{78}
{
    \zR{贾政才下衙门，正向贾琏问起 拿车之事。}
}
{
    \eR{[English source not present in the checked-in Bencraft/Joly files.]}
}
{
}

\Verse{79}
{
    \zR{贾琏道：“今儿叫人拿帖儿去，知县不在家。}
}
{
    \eR{[English source not present in the checked-in Bencraft/Joly files.]}
}
{
}

\Verse{80}
{
    \zR{他的门上说了：‘这是本 官不知道的，并无牌票出去拿车，都是那些混帐东西在外头撒野挤讹头。}
}
{
    \eR{[English source not present in the checked-in Bencraft/Joly files.]}
}
{
}

\Verse{81}
{
    \zR{既是老爷 府里的，我便立刻叫人去追办，包管明儿连车连东西一并送来。}
}
{
    \eR{[English source not present in the checked-in Bencraft/Joly files.]}
}
{
}

\Verse{82}
{
    \zR{如有半点差迟，再 行禀过本官，重重处治。}
}
{
    \eR{[English source not present in the checked-in Bencraft/Joly files.]}
}
{
}

\Verse{83}
{
    \zR{此刻本官不在家，求这里老爷看破些，可以不用本官知道 更好。’”}
}
{
    \eR{[English source not present in the checked-in Bencraft/Joly files.]}
}
{
}

\Verse{84}
{
    \zR{贾政道：“既无官票，到底是何等样人在那里作怪？”}
}
{
    \eR{[English source not present in the checked-in Bencraft/Joly files.]}
}
{
}

\Verse{85}
{
    \zR{贾琏道：“老爷 不知，外头都是这样。}
}
{
    \eR{[English source not present in the checked-in Bencraft/Joly files.]}
}
{
}

\Verse{86}
{
    \zR{想来明儿必定送来的。”}
}
{
    \eR{[English source not present in the checked-in Bencraft/Joly files.]}
}
{
}

\Verse{87}
{
    \zR{贾琏说完下来，宝玉上去见了。}
}
{
    \eR{[English source not present in the checked-in Bencraft/Joly files.]}
}
{
}

\Verse{88}
{
    \zR{贾 政问了几句，便叫他往老太太那里去。}
}
{
    \eR{[English source not present in the checked-in Bencraft/Joly files.]}
}
{
}

\Verse{89}
{
    \zR{贾琏因为昨夜叫空了家人，出来传唤，那起人都已伺候齐全。}
}
{
    \eR{[English source not present in the checked-in Bencraft/Joly files.]}
}
{
}

\Verse{90}
{
    \zR{贾琏骂了一顿， 叫大管家赖大：“将各行档的花名册子拿来，你去查点查点，写一张谕帖，叫那些 人知道。}
}
{
    \eR{[English source not present in the checked-in Bencraft/Joly files.]}
}
{
}

\Verse{91}
{
    \zR{若有并未告假，私自出去，传唤不到，贻误公事的，立刻给我打了撵出去！”}
}
{
    \eR{[English source not present in the checked-in Bencraft/Joly files.]}
}
{
}

\Verse{92}
{
    \zR{赖大连忙答应了几个“是”，出来吩咐了一回，家人各自留意。}
}
{
    \eR{[English source not present in the checked-in Bencraft/Joly files.]}
}
{
}

\Verse{93}
{
    \zR{过不几时，忽见有一个人，头上戴着毡帽，身上穿着一身青布衣裳，脚下穿着 一双撒鞋，走到门上，向众人作了个揖。}
}
{
    \eR{[English source not present in the checked-in Bencraft/Joly files.]}
}
{
}

\Verse{94}
{
    \zR{众人拿眼上上下下打量了他一番，便问他： “是那里来的？”}
}
{
    \eR{[English source not present in the checked-in Bencraft/Joly files.]}
}
{
}

\Verse{95}
{
    \zR{那人道：“我自南边甄府中来的。}
}
{
    \eR{[English source not present in the checked-in Bencraft/Joly files.]}
}
{
}

\Verse{96}
{
    \zR{并有家老爷手书一封，求这里 的爷们呈上尊老爷。”}
}
{
    \eR{[English source not present in the checked-in Bencraft/Joly files.]}
}
{
}

\Verse{97}
{
    \zR{众人听见他是甄府来的，才站起来让他坐下，道：“你乏了， 且坐坐。}
}
{
    \eR{[English source not present in the checked-in Bencraft/Joly files.]}
}
{
}

\Verse{98}
{
    \zR{我们给你回就是了。”}
}
{
    \eR{[English source not present in the checked-in Bencraft/Joly files.]}
}
{
}

\Verse{99}
{
    \zR{门上一面进来回明贾政，呈上来书。}
}
{
    \eR{[English source not present in the checked-in Bencraft/Joly files.]}
}
{
}

\Verse{100}
{
    \zR{贾政拆书看时， 上写着： 世交夙好，气谊素敦，遥仰？}
}
{
    \eR{[English source not present in the checked-in Bencraft/Joly files.]}
}
{
}

\Verse{101}
{
    \zR{帷，不胜依切。}
}
{
    \eR{[English source not present in the checked-in Bencraft/Joly files.]}
}
{
}

\Verse{102}
{
    \zR{弟因菲材获谴，自分万死难偿， 幸邀宽宥，待罪边隅。}
}
{
    \eR{[English source not present in the checked-in Bencraft/Joly files.]}
}
{
}

\Verse{103}
{
    \zR{迄今门户雕零，家人星散。}
}
{
    \eR{[English source not present in the checked-in Bencraft/Joly files.]}
}
{
}

\Verse{104}
{
    \zR{所有奴子包勇，向曾使用，虽无 奇技，人尚悫实。}
}
{
    \eR{[English source not present in the checked-in Bencraft/Joly files.]}
}
{
}

\Verse{105}
{
    \zR{倘使得备奔走，糊口有资，屋乌之爱，感佩无涯矣!专此奉达， 馀容再叙，不宣。}
}
{
    \eR{[English source not present in the checked-in Bencraft/Joly files.]}
}
{
}

\Verse{106}
{
    \zR{年家眷弟甄应嘉顿首。}
}
{
    \eR{[English source not present in the checked-in Bencraft/Joly files.]}
}
{
}

\Verse{107}
{
    \zR{贾政看完，笑道：“这里正因人多，甄家倒荐人来。}
}
{
    \eR{[English source not present in the checked-in Bencraft/Joly files.]}
}
{
}

\Verse{108}
{
    \zR{又不好却的。”}
}
{
    \eR{[English source not present in the checked-in Bencraft/Joly files.]}
}
{
}

\Verse{109}
{
    \zR{吩咐门上：“叫 他见我，且留他住下，因材使用便了。”}
}
{
    \eR{[English source not present in the checked-in Bencraft/Joly files.]}
}
{
}

\Verse{110}
{
    \zR{门上出去，带进人来，见贾政，便磕了三个头，起来道：“家老爷请老爷安。”}
}
{
    \eR{[English source not present in the checked-in Bencraft/Joly files.]}
}
{
}

\Verse{111}
{
    \zR{自己又打个千儿，说：“包勇请老爷安。”}
}
{
    \eR{[English source not present in the checked-in Bencraft/Joly files.]}
}
{
}

\Verse{112}
{
    \zR{贾政回问了甄老爷的好，便把他上下一 瞧。}
}
{
    \eR{[English source not present in the checked-in Bencraft/Joly files.]}
}
{
}

\Verse{113}
{
    \zR{但见包勇身长五尺有零，肩背宽肥，浓眉爆眼，磕额长髯，气色粗黑，垂着手 站着。}
}
{
    \eR{[English source not present in the checked-in Bencraft/Joly files.]}
}
{
}

\Verse{114}
{
    \zR{便问道：“你是向来在甄家的，还是住过几年的？”}
}
{
    \eR{[English source not present in the checked-in Bencraft/Joly files.]}
}
{
}

\Verse{115}
{
    \zR{包勇道：“小的向在甄 家的。”}
}
{
    \eR{[English source not present in the checked-in Bencraft/Joly files.]}
}
{
}

\Verse{116}
{
    \zR{贾政道：“你如今为什么要出来呢？”}
}
{
    \eR{[English source not present in the checked-in Bencraft/Joly files.]}
}
{
}

\Verse{117}
{
    \zR{包勇道：“小的原不肯出来，只是 家老爷再四叫小的出来，说别处你不肯去，这里老爷家里和在咱们自己家里一样 的，所以小的来的。”}
}
{
    \eR{[English source not present in the checked-in Bencraft/Joly files.]}
}
{
}

\Verse{118}
{
    \zR{贾政道：“你们老爷不该有这样事情，弄到这个田地。”}
}
{
    \eR{[English source not present in the checked-in Bencraft/Joly files.]}
}
{
}

\Verse{119}
{
    \zR{包 勇道：“小的本不敢说：我们老爷只是太好了，一味的真心待人，反倒招出事来。”}
}
{
    \eR{[English source not present in the checked-in Bencraft/Joly files.]}
}
{
}

\Verse{120}
{
    \zR{贾政道：“真心是最好的了。”}
}
{
    \eR{[English source not present in the checked-in Bencraft/Joly files.]}
}
{
}

\Verse{121}
{
    \zR{包勇道：“因为太真了，人人都不喜欢，讨人厌烦 是有的。”}
}
{
    \eR{[English source not present in the checked-in Bencraft/Joly files.]}
}
{
}

\Verse{122}
{
    \zR{贾政笑了一笑道：“既这样，皇天自然不负他的。”}
}
{
    \eR{[English source not present in the checked-in Bencraft/Joly files.]}
}
{
}

\Verse{123}
{
    \zR{包勇还要说时，贾 政又问道：“我听见说你们家的哥儿不是也叫宝玉么？”}
}
{
    \eR{[English source not present in the checked-in Bencraft/Joly files.]}
}
{
}

\Verse{124}
{
    \zR{包勇道：“是。”}
}
{
    \eR{[English source not present in the checked-in Bencraft/Joly files.]}
}
{
}

\Verse{125}
{
    \zR{贾政道： “他还肯向上巴结么？”}
}
{
    \eR{[English source not present in the checked-in Bencraft/Joly files.]}
}
{
}

\Verse{126}
{
    \zR{包勇道：“老爷若问我们哥儿，倒是一段奇事。}
}
{
    \eR{[English source not present in the checked-in Bencraft/Joly files.]}
}
{
}

\Verse{127}
{
    \zR{哥儿的脾 气也和我家老爷一个样子，也是一味的诚实，从小儿只爱和那些姐妹们在一处玩。}
}
{
    \eR{[English source not present in the checked-in Bencraft/Joly files.]}
}
{
}

\Verse{128}
{
    \zR{老爷太太也狠打过几次，他只是不改。}
}
{
    \eR{[English source not present in the checked-in Bencraft/Joly files.]}
}
{
}

\Verse{129}
{
    \zR{那一年太太进京的时候儿，哥儿大病了一场， 已经死了半日，把老爷几乎急死，装裹都预备了。}
}
{
    \eR{[English source not present in the checked-in Bencraft/Joly files.]}
}
{
}

\Verse{130}
{
    \zR{幸喜后来好了，嘴里说道：走到 一座牌楼那里，见了一个姑娘，领着他到了一座庙里，见了好些柜子，里头见了好 些册子。}
}
{
    \eR{[English source not present in the checked-in Bencraft/Joly files.]}
}
{
}

\Verse{131}
{
    \zR{又到屋里，见了无数女子，说是都变了鬼怪似的，也有变做骷髅儿的。}
}
{
    \eR{[English source not present in the checked-in Bencraft/Joly files.]}
}
{
}

\Verse{132}
{
    \zR{他 吓急了，就哭喊起来。}
}
{
    \eR{[English source not present in the checked-in Bencraft/Joly files.]}
}
{
}

\Verse{133}
{
    \zR{老爷知他醒过来了，连忙调治，渐渐的好了。}
}
{
    \eR{[English source not present in the checked-in Bencraft/Joly files.]}
}
{
}

\Verse{134}
{
    \zR{老爷仍叫他在 姐妹们一处玩去，他竟改了脾气了：好着时候的玩意儿一概都不要了，惟有念书为 事。}
}
{
    \eR{[English source not present in the checked-in Bencraft/Joly files.]}
}
{
}

\Verse{135}
{
    \zR{就有什么人来引诱他，他也全不动心。}
}
{
    \eR{[English source not present in the checked-in Bencraft/Joly files.]}
}
{
}

\Verse{136}
{
    \zR{如今渐渐的能够帮着老爷料理些家务 了。”}
}
{
    \eR{[English source not present in the checked-in Bencraft/Joly files.]}
}
{
}

\Verse{137}
{
    \zR{贾政默然想了一回，道：“你去歇歇去罢。}
}
{
    \eR{[English source not present in the checked-in Bencraft/Joly files.]}
}
{
}

\Verse{138}
{
    \zR{等这里用着你时，自然派你一个 行次儿。”}
}
{
    \eR{[English source not present in the checked-in Bencraft/Joly files.]}
}
{
}

\Verse{139}
{
    \zR{包勇答应着，退下来，跟着这里人出去歇息不提。}
}
{
    \eR{[English source not present in the checked-in Bencraft/Joly files.]}
}
{
}

\Verse{140}
{
    \zR{一日贾政早起，刚要上衙门，看见门上那些人在那里交头接耳，好像要使贾政 知道的似的，又不好明回，只管咕咕唧唧的说话。}
}
{
    \eR{[English source not present in the checked-in Bencraft/Joly files.]}
}
{
}

\Verse{141}
{
    \zR{贾政叫上来问道：“你们有什么 事这么鬼鬼祟祟的？”}
}
{
    \eR{[English source not present in the checked-in Bencraft/Joly files.]}
}
{
}

\Verse{142}
{
    \zR{门上的人回道：“奴才们不敢说。”}
}
{
    \eR{[English source not present in the checked-in Bencraft/Joly files.]}
}
{
}

\Verse{143}
{
    \zR{贾政道：“有什么事不 敢说的？”}
}
{
    \eR{[English source not present in the checked-in Bencraft/Joly files.]}
}
{
}

\Verse{144}
{
    \zR{门上的人道：“奴才今儿起来，开门出去，见门上贴着一张白纸，上写 着许多不成事体的字。”}
}
{
    \eR{[English source not present in the checked-in Bencraft/Joly files.]}
}
{
}

\Verse{145}
{
    \zR{贾政道：“那里有这样的事!写的是什么？”}
}
{
    \eR{[English source not present in the checked-in Bencraft/Joly files.]}
}
{
}

\Verse{146}
{
    \zR{门上的人道： “是水月庵里的腌？}
}
{
    \eR{[English source not present in the checked-in Bencraft/Joly files.]}
}
{
}

\Verse{147}
{
    \zR{话。”}
}
{
    \eR{[English source not present in the checked-in Bencraft/Joly files.]}
}
{
}

\Verse{148}
{
    \zR{贾政道：“拿给我瞧。”}
}
{
    \eR{[English source not present in the checked-in Bencraft/Joly files.]}
}
{
}

\Verse{149}
{
    \zR{门上的人道：“奴才本要揭下 来，谁知他贴的结实，揭不下来，只得一面抄，一面洗。}
}
{
    \eR{[English source not present in the checked-in Bencraft/Joly files.]}
}
{
}

\Verse{150}
{
    \zR{刚才李德揭了一张给奴才 瞧，就是那门上贴的话。}
}
{
    \eR{[English source not present in the checked-in Bencraft/Joly files.]}
}
{
}

\Verse{151}
{
    \zR{奴才们不敢隐瞒。”}
}
{
    \eR{[English source not present in the checked-in Bencraft/Joly files.]}
}
{
}

\Verse{152}
{
    \zR{说着，呈上那帖儿。}
}
{
    \eR{[English source not present in the checked-in Bencraft/Joly files.]}
}
{
}

\Verse{153}
{
    \zR{贾政接来看时， 上面写着： 西贝草斤年纪轻，水月庵里管尼僧。}
}
{
    \eR{[English source not present in the checked-in Bencraft/Joly files.]}
}
{
}

\Verse{154}
{
    \zR{一个男人多少女，窝娼聚赌是陶情。}
}
{
    \eR{[English source not present in the checked-in Bencraft/Joly files.]}
}
{
}

\Verse{155}
{
    \zR{不肖子弟来办事，荣国府内好声名。}
}
{
    \eR{[English source not present in the checked-in Bencraft/Joly files.]}
}
{
}

\Verse{156}
{
    \zR{贾政看了，气的头昏目晕，赶着叫门上的人不许声张，悄悄叫人往宁荣两府靠 近的夹道子墙壁上再去找寻。}
}
{
    \eR{[English source not present in the checked-in Bencraft/Joly files.]}
}
{
}

\Verse{157}
{
    \zR{随即叫人去唤贾琏出来。}
}
{
    \eR{[English source not present in the checked-in Bencraft/Joly files.]}
}
{
}

\Verse{158}
{
    \zR{贾琏即忙赶至。}
}
{
    \eR{[English source not present in the checked-in Bencraft/Joly files.]}
}
{
}

\Verse{159}
{
    \zR{贾政忙问道： “水月庵中寄居的那些女尼女道，向来你也查考查考过没有？”}
}
{
    \eR{[English source not present in the checked-in Bencraft/Joly files.]}
}
{
}

\Verse{160}
{
    \zR{贾琏道：“没有， 一向都是芹儿在那里照管。”}
}
{
    \eR{[English source not present in the checked-in Bencraft/Joly files.]}
}
{
}

\Verse{161}
{
    \zR{贾政道：“你知道芹儿照管得来照管不来？”}
}
{
    \eR{[English source not present in the checked-in Bencraft/Joly files.]}
}
{
}

\Verse{162}
{
    \zR{贾琏道： “老爷既这么说，想来芹儿必有不妥当的地方儿。”}
}
{
    \eR{[English source not present in the checked-in Bencraft/Joly files.]}
}
{
}

\Verse{163}
{
    \zR{贾政叹道：“你瞧瞧这个帖儿 写的是什么！”}
}
{
    \eR{[English source not present in the checked-in Bencraft/Joly files.]}
}
{
}

\Verse{164}
{
    \zR{贾琏一看道：“有这样事么！”}
}
{
    \eR{[English source not present in the checked-in Bencraft/Joly files.]}
}
{
}

\Verse{165}
{
    \zR{正说着，只见贾蓉走来，拿着一封 书子，写着“二老爷密启”。}
}
{
    \eR{[English source not present in the checked-in Bencraft/Joly files.]}
}
{
}

\Verse{166}
{
    \zR{打开看时，也是无头榜一张，与门上所贴的话相同。}
}
{
    \eR{[English source not present in the checked-in Bencraft/Joly files.]}
}
{
}

\Verse{167}
{
    \zR{贾政道：“快叫赖大带了三四辆车到水月庵里去，把那些女尼姑女道士一齐拉回来。}
}
{
    \eR{[English source not present in the checked-in Bencraft/Joly files.]}
}
{
}

\Verse{168}
{
    \zR{不许泄漏，只说里头传唤。”}
}
{
    \eR{[English source not present in the checked-in Bencraft/Joly files.]}
}
{
}

\Verse{169}
{
    \zR{赖大领命去了。}
}
{
    \eR{[English source not present in the checked-in Bencraft/Joly files.]}
}
{
}

\Verse{170}
{
    \zR{且说水月庵中小女尼女道士等，初到庵中，沙弥与道士原系老尼收管，日间教 他些经忏。}
}
{
    \eR{[English source not present in the checked-in Bencraft/Joly files.]}
}
{
}

\Verse{171}
{
    \zR{以后元妃不用，也便习学得懒惰了。}
}
{
    \eR{[English source not present in the checked-in Bencraft/Joly files.]}
}
{
}

\Verse{172}
{
    \zR{那些女孩子们年纪渐渐的大了，都 也有些知觉了。}
}
{
    \eR{[English source not present in the checked-in Bencraft/Joly files.]}
}
{
}

\Verse{173}
{
    \zR{更兼贾芹也是风流人物，打量芳官等出家，只是小孩子性儿，便去 招惹他们。}
}
{
    \eR{[English source not present in the checked-in Bencraft/Joly files.]}
}
{
}

\Verse{174}
{
    \zR{那知芳官竟是真心，不能上手，便把这心肠移到女尼女道士身上。}
}
{
    \eR{[English source not present in the checked-in Bencraft/Joly files.]}
}
{
}

\Verse{175}
{
    \zR{因那 小沙弥中有个名叫沁香的，和女道士中有个叫做鹤仙的，长的都甚妖娆，贾芹便和 这两个人勾搭上了，闲时便学些丝弦，唱个曲儿。}
}
{
    \eR{[English source not present in the checked-in Bencraft/Joly files.]}
}
{
}

\Verse{176}
{
    \zR{那时正当十月中旬，贾芹给庵中那些人领了月例银子，便想起法儿来，告诉众 人道：“我为你们领月钱，不能进城，又只得在这里歇着，怪冷的。}
}
{
    \eR{[English source not present in the checked-in Bencraft/Joly files.]}
}
{
}

\Verse{177}
{
    \zR{怎么样？}
}
{
    \eR{[English source not present in the checked-in Bencraft/Joly files.]}
}
{
}

\Verse{178}
{
    \zR{我今 儿带些果子酒，大家吃着乐一夜好不好？”}
}
{
    \eR{[English source not present in the checked-in Bencraft/Joly files.]}
}
{
}

\Verse{179}
{
    \zR{那些女孩子都高兴，便摆起桌子，连本 庵的女尼也叫了来。}
}
{
    \eR{[English source not present in the checked-in Bencraft/Joly files.]}
}
{
}

\Verse{180}
{
    \zR{惟有芳官不来。}
}
{
    \eR{[English source not present in the checked-in Bencraft/Joly files.]}
}
{
}

\Verse{181}
{
    \zR{贾芹喝了几杯，便说道要行令。}
}
{
    \eR{[English source not present in the checked-in Bencraft/Joly files.]}
}
{
}

\Verse{182}
{
    \zR{沁香等道：“我 们都不会，倒不如？}
}
{
    \eR{[English source not present in the checked-in Bencraft/Joly files.]}
}
{
}

\Verse{183}
{
    \zR{拳罢。}
}
{
    \eR{[English source not present in the checked-in Bencraft/Joly files.]}
}
{
}

\Verse{184}
{
    \zR{谁输了喝一钟，岂不爽快？”}
}
{
    \eR{[English source not present in the checked-in Bencraft/Joly files.]}
}
{
}

\Verse{185}
{
    \zR{本庵的女尼道：“这天刚 过晌午，混嚷混喝的不像，且先喝几钟，爱散的先散去。}
}
{
    \eR{[English source not present in the checked-in Bencraft/Joly files.]}
}
{
}

\Verse{186}
{
    \zR{谁爱陪芹大爷的，回来晚 上尽子喝去，我也不管。”}
}
{
    \eR{[English source not present in the checked-in Bencraft/Joly files.]}
}
{
}

\Verse{187}
{
    \zR{正说着，只见道婆急忙进来说：“快散了罢!府里赖大 爷来了。”}
}
{
    \eR{[English source not present in the checked-in Bencraft/Joly files.]}
}
{
}

\Verse{188}
{
    \zR{众女尼忙乱收拾，便叫贾芹躲开。}
}
{
    \eR{[English source not present in the checked-in Bencraft/Joly files.]}
}
{
}

\Verse{189}
{
    \zR{贾芹因多喝了几杯，便道：“我是送 月钱来的，怕什么？”}
}
{
    \eR{[English source not present in the checked-in Bencraft/Joly files.]}
}
{
}

\Verse{190}
{
    \zR{话犹未完，已见赖大进来，见这般样子，心里大怒。}
}
{
    \eR{[English source not present in the checked-in Bencraft/Joly files.]}
}
{
}

\Verse{191}
{
    \zR{为的是 贾政吩咐不许声张，只得含糊装笑道：“芹大爷也在这里呢么？”}
}
{
    \eR{[English source not present in the checked-in Bencraft/Joly files.]}
}
{
}

\Verse{192}
{
    \zR{贾芹连忙站起来 道：“赖大爷，你来作什么？”}
}
{
    \eR{[English source not present in the checked-in Bencraft/Joly files.]}
}
{
}

\Verse{193}
{
    \zR{赖大说：“大爷在这里更好。}
}
{
    \eR{[English source not present in the checked-in Bencraft/Joly files.]}
}
{
}

\Verse{194}
{
    \zR{快快叫沙弥道士收拾 上车进城，宫里传呢。”}
}
{
    \eR{[English source not present in the checked-in Bencraft/Joly files.]}
}
{
}

\Verse{195}
{
    \zR{贾芹等不知原故，还要细问。}
}
{
    \eR{[English source not present in the checked-in Bencraft/Joly files.]}
}
{
}

\Verse{196}
{
    \zR{赖大说：“天已不早了，快 快的好赶进城。”}
}
{
    \eR{[English source not present in the checked-in Bencraft/Joly files.]}
}
{
}

\Verse{197}
{
    \zR{众女孩子只得一齐上车。}
}
{
    \eR{[English source not present in the checked-in Bencraft/Joly files.]}
}
{
}

\Verse{198}
{
    \zR{赖大骑着大走骡，押着赶进城，不提。}
}
{
    \eR{[English source not present in the checked-in Bencraft/Joly files.]}
}
{
}

\Verse{199}
{
    \zR{却说贾政知道这事，气的衙门也不能上了，独坐在内书房叹气。}
}
{
    \eR{[English source not present in the checked-in Bencraft/Joly files.]}
}
{
}

\Verse{200}
{
    \zR{贾琏也不敢走 开。}
}
{
    \eR{[English source not present in the checked-in Bencraft/Joly files.]}
}
{
}

\Verse{201}
{
    \zR{忽见门上的进来禀道：“衙门里今夜该班是张老爷。}
}
{
    \eR{[English source not present in the checked-in Bencraft/Joly files.]}
}
{
}

\Verse{202}
{
    \zR{因张老爷病了，有知会来 请老爷补一班。”}
}
{
    \eR{[English source not present in the checked-in Bencraft/Joly files.]}
}
{
}

\Verse{203}
{
    \zR{贾政正等赖大回来要办贾芹，此时又要该班，心里纳闷，也不言 语。}
}
{
    \eR{[English source not present in the checked-in Bencraft/Joly files.]}
}
{
}

\Verse{204}
{
    \zR{贾琏走上去说道：“赖大是饭后出去的，水月庵离城二十来里，就赶进城也得 二更天。}
}
{
    \eR{[English source not present in the checked-in Bencraft/Joly files.]}
}
{
}

\Verse{205}
{
    \zR{今日又是老爷的帮班，请老爷只管去。}
}
{
    \eR{[English source not present in the checked-in Bencraft/Joly files.]}
}
{
}

\Verse{206}
{
    \zR{赖大来了，叫他押着，也别声张， 等明儿老爷回来再发落。}
}
{
    \eR{[English source not present in the checked-in Bencraft/Joly files.]}
}
{
}

\Verse{207}
{
    \zR{倘或芹儿来了，也不用说明，看他明儿见了老爷怎么样 说。”}
}
{
    \eR{[English source not present in the checked-in Bencraft/Joly files.]}
}
{
}

\Verse{208}
{
    \zR{贾政听来有理，只得上班去了。}
}
{
    \eR{[English source not present in the checked-in Bencraft/Joly files.]}
}
{
}

\Verse{209}
{
    \zR{贾琏抽空才要回到自己房中，一面走着，心 里抱怨凤姐出的主意，欲要埋怨，因他病着，只得隐忍，慢慢的走着。}
}
{
    \eR{[English source not present in the checked-in Bencraft/Joly files.]}
}
{
}

\Verse{210}
{
    \zR{且说那些下人，一人传十，传到里头，先是平儿知道，即忙告诉凤姐。}
}
{
    \eR{[English source not present in the checked-in Bencraft/Joly files.]}
}
{
}

\Verse{211}
{
    \zR{凤姐因 那一夜不好，恹恹的总没精神，正是惦记铁槛寺的事情。}
}
{
    \eR{[English source not present in the checked-in Bencraft/Joly files.]}
}
{
}

\Verse{212}
{
    \zR{听见“外头贴了匿名揭帖” 的一句话，吓了一跳，忙问：“贴的是什么？”}
}
{
    \eR{[English source not present in the checked-in Bencraft/Joly files.]}
}
{
}

\Verse{213}
{
    \zR{平儿随口答应，不留神，就错说了， 道：“没要紧，是馒头庵里的事情。”}
}
{
    \eR{[English source not present in the checked-in Bencraft/Joly files.]}
}
{
}

\Verse{214}
{
    \zR{凤姐本是心虚，听见“馒头庵的事情”，这 一唬直唬怔了，一句话没说出来，急火上攻，眼前发晕，咳嗽了一阵便歪倒了，两 只眼却只是发怔。}
}
{
    \eR{[English source not present in the checked-in Bencraft/Joly files.]}
}
{
}

\Verse{215}
{
    \zR{平儿慌了，说道：“水月庵里，不过是女沙弥女道士的事，奶奶 着什么急呢？”}
}
{
    \eR{[English source not present in the checked-in Bencraft/Joly files.]}
}
{
}

\Verse{216}
{
    \zR{凤姐听是水月庵，才定了定神，道：“嗳!糊涂东西!到底是水月庵， 是馒头庵呢？”}
}
{
    \eR{[English source not present in the checked-in Bencraft/Joly files.]}
}
{
}

\Verse{217}
{
    \zR{平儿道：“是我头里错听了馒头庵，后来听见不是馒头庵，是水月 庵。}
}
{
    \eR{[English source not present in the checked-in Bencraft/Joly files.]}
}
{
}

\Verse{218}
{
    \zR{我刚才也就说溜了嘴，说成馒头庵了。”}
}
{
    \eR{[English source not present in the checked-in Bencraft/Joly files.]}
}
{
}

\Verse{219}
{
    \zR{凤姐道：“我就知道是水月庵。}
}
{
    \eR{[English source not present in the checked-in Bencraft/Joly files.]}
}
{
}

\Verse{220}
{
    \zR{那馒 头庵与我什么相干。}
}
{
    \eR{[English source not present in the checked-in Bencraft/Joly files.]}
}
{
}

\Verse{221}
{
    \zR{原是这水月庵是我叫芹儿管的，大约刻扣了月钱。”}
}
{
    \eR{[English source not present in the checked-in Bencraft/Joly files.]}
}
{
}

\Verse{222}
{
    \zR{平儿道： “我听着不像月钱的事，还有些腌？}
}
{
    \eR{[English source not present in the checked-in Bencraft/Joly files.]}
}
{
}

\Verse{223}
{
    \zR{话呢。”}
}
{
    \eR{[English source not present in the checked-in Bencraft/Joly files.]}
}
{
}

\Verse{224}
{
    \zR{凤姐道：“我更不管那个。}
}
{
    \eR{[English source not present in the checked-in Bencraft/Joly files.]}
}
{
}

\Verse{225}
{
    \zR{你二爷那 里去了？”}
}
{
    \eR{[English source not present in the checked-in Bencraft/Joly files.]}
}
{
}

\Verse{226}
{
    \zR{平儿说：“听见老爷生气，他不敢走开。}
}
{
    \eR{[English source not present in the checked-in Bencraft/Joly files.]}
}
{
}

\Verse{227}
{
    \zR{我听见事情不好，我吩咐这些 人不许吵嚷，不知太太们知道了没有。}
}
{
    \eR{[English source not present in the checked-in Bencraft/Joly files.]}
}
{
}

\Verse{228}
{
    \zR{就听见说，老爷叫赖大拿这些女孩子去了。}
}
{
    \eR{[English source not present in the checked-in Bencraft/Joly files.]}
}
{
}

\Verse{229}
{
    \zR{且叫人前头打听打听。}
}
{
    \eR{[English source not present in the checked-in Bencraft/Joly files.]}
}
{
}

\Verse{230}
{
    \zR{奶奶现在病着，依我竟先别管他们的闲事。”}
}
{
    \eR{[English source not present in the checked-in Bencraft/Joly files.]}
}
{
}

\Verse{231}
{
    \zR{正说着，只见 贾琏进来。}
}
{
    \eR{[English source not present in the checked-in Bencraft/Joly files.]}
}
{
}

\Verse{232}
{
    \zR{凤姐欲待问他，见贾琏一脸怒气，暂且装作不知。}
}
{
    \eR{[English source not present in the checked-in Bencraft/Joly files.]}
}
{
}

\Verse{233}
{
    \zR{贾琏没吃完饭，旺儿 来说：“外头请爷呢，赖大回来了。”}
}
{
    \eR{[English source not present in the checked-in Bencraft/Joly files.]}
}
{
}

\Verse{234}
{
    \zR{贾琏道：“芹儿来了没有？”}
}
{
    \eR{[English source not present in the checked-in Bencraft/Joly files.]}
}
{
}

\Verse{235}
{
    \zR{旺儿道：“也 来了。”}
}
{
    \eR{[English source not present in the checked-in Bencraft/Joly files.]}
}
{
}

\Verse{236}
{
    \zR{贾琏便道：“你去告诉赖大说：老爷上班儿去了，把这些个女孩子暂且收 在园里，明日等老爷回来，送进宫去。}
}
{
    \eR{[English source not present in the checked-in Bencraft/Joly files.]}
}
{
}

\Verse{237}
{
    \zR{只叫芹儿在内书房等着我。”}
}
{
    \eR{[English source not present in the checked-in Bencraft/Joly files.]}
}
{
}

\Verse{238}
{
    \zR{旺儿去了。}
}
{
    \eR{[English source not present in the checked-in Bencraft/Joly files.]}
}
{
}

\Verse{239}
{
    \zR{贾芹走进书房，只见那些下人指指戳戳不知说什么，看起这个样儿来，不像宫 里要人。}
}
{
    \eR{[English source not present in the checked-in Bencraft/Joly files.]}
}
{
}

\Verse{240}
{
    \zR{想着问人，又问不出来。}
}
{
    \eR{[English source not present in the checked-in Bencraft/Joly files.]}
}
{
}

\Verse{241}
{
    \zR{正在心里疑惑，只见贾琏走出来，贾芹便请了安， 垂手侍立，说道：“不知道娘娘宫里即刻传那些孩子们做什么？}
}
{
    \eR{[English source not present in the checked-in Bencraft/Joly files.]}
}
{
}

\Verse{242}
{
    \zR{叫侄儿好赶。}
}
{
    \eR{[English source not present in the checked-in Bencraft/Joly files.]}
}
{
}

\Verse{243}
{
    \zR{幸喜 侄儿今儿送月钱去，还没有走，便同着赖大来了。}
}
{
    \eR{[English source not present in the checked-in Bencraft/Joly files.]}
}
{
}

\Verse{244}
{
    \zR{二叔想来是知道的。”}
}
{
    \eR{[English source not present in the checked-in Bencraft/Joly files.]}
}
{
}

\Verse{245}
{
    \zR{贾琏道： “我知道什么？}
}
{
    \eR{[English source not present in the checked-in Bencraft/Joly files.]}
}
{
}

\Verse{246}
{
    \zR{你才是明白的呢！”}
}
{
    \eR{[English source not present in the checked-in Bencraft/Joly files.]}
}
{
}

\Verse{247}
{
    \zR{贾芹摸不着头脑儿，也不敢再问。}
}
{
    \eR{[English source not present in the checked-in Bencraft/Joly files.]}
}
{
}

\Verse{248}
{
    \zR{贾琏道：“你 干的好事啊!把老爷都气坏了！”}
}
{
    \eR{[English source not present in the checked-in Bencraft/Joly files.]}
}
{
}

\Verse{249}
{
    \zR{贾芹道：“侄儿没有干什么。}
}
{
    \eR{[English source not present in the checked-in Bencraft/Joly files.]}
}
{
}

\Verse{250}
{
    \zR{庵里月钱是月月给 的，孩子们经忏是不忘的。”}
}
{
    \eR{[English source not present in the checked-in Bencraft/Joly files.]}
}
{
}

\Verse{251}
{
    \zR{贾琏见他不知，又是平素常在一处玩笑的，便叹口气 道：“打嘴的东西，你各自去瞧瞧罢。”}
}
{
    \eR{[English source not present in the checked-in Bencraft/Joly files.]}
}
{
}

\Verse{252}
{
    \zR{便从靴掖儿里头拿出那个揭帖来，扔与他 瞧。}
}
{
    \eR{[English source not present in the checked-in Bencraft/Joly files.]}
}
{
}

\Verse{253}
{
    \zR{贾芹拾来一看，吓得面如土色，说道：“这是谁干的!我并没得罪人，为什么 这么坑我？}
}
{
    \eR{[English source not present in the checked-in Bencraft/Joly files.]}
}
{
}

\Verse{254}
{
    \zR{我一月送钱去，只走一趟，并没有这些事。}
}
{
    \eR{[English source not present in the checked-in Bencraft/Joly files.]}
}
{
}

\Verse{255}
{
    \zR{若是老爷回来，打着问我， 侄儿就屈死了!我母亲知道，更要打死。”}
}
{
    \eR{[English source not present in the checked-in Bencraft/Joly files.]}
}
{
}

\Verse{256}
{
    \zR{说着，见没人在旁边，便跪下央及道： “好叔叔，救我一救儿罢！”}
}
{
    \eR{[English source not present in the checked-in Bencraft/Joly files.]}
}
{
}

\Verse{257}
{
    \zR{说着，只管磕头，满眼流泪。}
}
{
    \eR{[English source not present in the checked-in Bencraft/Joly files.]}
}
{
}

\Verse{258}
{
    \zR{贾琏想道：“老爷最恼 这些，要是问准了有这些事，这场气也不小，闹出去也不好听。}
}
{
    \eR{[English source not present in the checked-in Bencraft/Joly files.]}
}
{
}

\Verse{259}
{
    \zR{又长那个贴帖儿的 人的志气了，将来咱们的事多着呢。}
}
{
    \eR{[English source not present in the checked-in Bencraft/Joly files.]}
}
{
}

\Verse{260}
{
    \zR{倒不如趁着老爷上班儿，和赖大商量着，要混 过去，就可以没事了。}
}
{
    \eR{[English source not present in the checked-in Bencraft/Joly files.]}
}
{
}

\Verse{261}
{
    \zR{现在没有对证。”}
}
{
    \eR{[English source not present in the checked-in Bencraft/Joly files.]}
}
{
}

\Verse{262}
{
    \zR{想定主意，便说：“你别瞒我。}
}
{
    \eR{[English source not present in the checked-in Bencraft/Joly files.]}
}
{
}

\Verse{263}
{
    \zR{你干的鬼 儿，你打量我都不知道呢。}
}
{
    \eR{[English source not present in the checked-in Bencraft/Joly files.]}
}
{
}

\Verse{264}
{
    \zR{若要完事，除非是老爷打着问你，你只一口咬定没有才 好。}
}
{
    \eR{[English source not present in the checked-in Bencraft/Joly files.]}
}
{
}

\Verse{265}
{
    \zR{没脸的东西!起去罢！”}
}
{
    \eR{[English source not present in the checked-in Bencraft/Joly files.]}
}
{
}

\Verse{266}
{
    \zR{叫人去叫赖大。}
}
{
    \eR{[English source not present in the checked-in Bencraft/Joly files.]}
}
{
}

\Verse{267}
{
    \zR{不多时，赖大来了，贾琏便和他商量。}
}
{
    \eR{[English source not present in the checked-in Bencraft/Joly files.]}
}
{
}

\Verse{268}
{
    \zR{赖大说：“这芹大爷本来闹的不像了。}
}
{
    \eR{[English source not present in the checked-in Bencraft/Joly files.]}
}
{
}

\Verse{269}
{
    \zR{奴才今儿到庵里的时候，他们正在那里喝酒呢。}
}
{
    \eR{[English source not present in the checked-in Bencraft/Joly files.]}
}
{
}

\Verse{270}
{
    \zR{帖儿上的话一定是有的。”}
}
{
    \eR{[English source not present in the checked-in Bencraft/Joly files.]}
}
{
}

\Verse{271}
{
    \zR{贾琏道： “芹儿，你听!赖大还赖你不成？”}
}
{
    \eR{[English source not present in the checked-in Bencraft/Joly files.]}
}
{
}

\Verse{272}
{
    \zR{贾芹此时红涨了脸，一句也不敢言语。}
}
{
    \eR{[English source not present in the checked-in Bencraft/Joly files.]}
}
{
}

\Verse{273}
{
    \zR{还是贾 琏拉着赖大，央他：“护庇护庇罢，只说芹哥儿是在家里找了来的。}
}
{
    \eR{[English source not present in the checked-in Bencraft/Joly files.]}
}
{
}

\Verse{274}
{
    \zR{你带了他去， 只说没有见我。}
}
{
    \eR{[English source not present in the checked-in Bencraft/Joly files.]}
}
{
}

\Verse{275}
{
    \zR{明日你求老爷，也不用问那些女孩子了，竟是叫了媒人来，领了去， 一卖完事。}
}
{
    \eR{[English source not present in the checked-in Bencraft/Joly files.]}
}
{
}

\Verse{276}
{
    \zR{果然娘娘再要的时候儿，咱们再买。”}
}
{
    \eR{[English source not present in the checked-in Bencraft/Joly files.]}
}
{
}

\Verse{277}
{
    \zR{赖大想来，闹也无益，且名声不 好，也就应了。}
}
{
    \eR{[English source not present in the checked-in Bencraft/Joly files.]}
}
{
}

\Verse{278}
{
    \zR{贾琏叫贾芹：“跟了赖大爷去罢!听着他教你，你就跟着他。”}
}
{
    \eR{[English source not present in the checked-in Bencraft/Joly files.]}
}
{
}

\Verse{279}
{
    \zR{说 罢，贾芹又磕了一个头，跟着赖大出去。}
}
{
    \eR{[English source not present in the checked-in Bencraft/Joly files.]}
}
{
}

\Verse{280}
{
    \zR{到了没人的地方儿，又给赖大磕头。}
}
{
    \eR{[English source not present in the checked-in Bencraft/Joly files.]}
}
{
}

\Verse{281}
{
    \zR{赖大 说：“我的小爷，你太闹的不像了。}
}
{
    \eR{[English source not present in the checked-in Bencraft/Joly files.]}
}
{
}

\Verse{282}
{
    \zR{不知得罪了谁，闹出这个乱儿来，你想想，谁 和你不对罢？”}
}
{
    \eR{[English source not present in the checked-in Bencraft/Joly files.]}
}
{
}

\Verse{283}
{
    \zR{贾芹想了一会子，并无不对的人，只得无精打彩，跟着赖大走回。}
}
{
    \eR{[English source not present in the checked-in Bencraft/Joly files.]}
}
{
}

\Verse{284}
{
    \zR{未知如何抵赖，且听下回分解。}
}
{
    \eR{[English source not present in the checked-in Bencraft/Joly files.]}
}
{
}

\Verse{285}
{
    \zR{(本章完)}
}
{
    \eR{[English source not present in the checked-in Bencraft/Joly files.]}
}
{
}
